% Options for packages loaded elsewhere
\PassOptionsToPackage{unicode}{hyperref}
\PassOptionsToPackage{hyphens}{url}
\PassOptionsToPackage{dvipsnames,svgnames,x11names}{xcolor}
%
\documentclass[
  10pt,
]{article}
\usepackage{amsmath,amssymb}
\usepackage{iftex}
\ifPDFTeX
  \usepackage[T1]{fontenc}
  \usepackage[utf8]{inputenc}
  \usepackage{textcomp} % provide euro and other symbols
\else % if luatex or xetex
  \usepackage{unicode-math} % this also loads fontspec
  \defaultfontfeatures{Scale=MatchLowercase}
  \defaultfontfeatures[\rmfamily]{Ligatures=TeX,Scale=1}
\fi
\usepackage{lmodern}
\ifPDFTeX\else
  % xetex/luatex font selection
  \setmainfont[]{DejaVu Serif}
  \setsansfont[]{DejaVu Sans}
  \setmonofont[]{DejaVu Sans Mono}
  \setmathfont[]{Latin Modern Math}
\fi
% Use upquote if available, for straight quotes in verbatim environments
\IfFileExists{upquote.sty}{\usepackage{upquote}}{}
\IfFileExists{microtype.sty}{% use microtype if available
  \usepackage[]{microtype}
  \UseMicrotypeSet[protrusion]{basicmath} % disable protrusion for tt fonts
}{}
\makeatletter
\@ifundefined{KOMAClassName}{% if non-KOMA class
  \IfFileExists{parskip.sty}{%
    \usepackage{parskip}
  }{% else
    \setlength{\parindent}{0pt}
    \setlength{\parskip}{6pt plus 2pt minus 1pt}}
}{% if KOMA class
  \KOMAoptions{parskip=half}}
\makeatother
\usepackage{xcolor}
\usepackage[margin=2.0cm]{geometry}
\usepackage{longtable,booktabs,array}
\usepackage{calc} % for calculating minipage widths
% Correct order of tables after \paragraph or \subparagraph
\usepackage{etoolbox}
\makeatletter
\patchcmd\longtable{\par}{\if@noskipsec\mbox{}\fi\par}{}{}
\makeatother
% Allow footnotes in longtable head/foot
\IfFileExists{footnotehyper.sty}{\usepackage{footnotehyper}}{\usepackage{footnote}}
\makesavenoteenv{longtable}
\setlength{\emergencystretch}{3em} % prevent overfull lines
\providecommand{\tightlist}{%
  \setlength{\itemsep}{0pt}\setlength{\parskip}{0pt}}
\setcounter{secnumdepth}{-\maxdimen} % remove section numbering
% doc_header.tex -- LaTeX preamble additions used by build_docs.sh for every
% Markdown -> LaTeX -> PDF conversion (the documents of the folder Revision_SM_Higgs).
\usepackage{etoolbox}
% longtable may not be loaded by pandoc when a document has no table
\usepackage{longtable}
% bold vectors: \bm -> \symbf under unicode-math (XeLaTeX), bm package otherwise
\makeatletter
\@ifpackageloaded{unicode-math}{\providecommand{\bm}[1]{\symbf{#1}}}{\usepackage{bm}}
\makeatother
% smaller type in long tables so that wide audit tables fit the text width;
% inside a table a line break is allowed after an underscore of plain text
% (pandoc writes it as \_), so that long identifiers wrap inside their column
% rows are set a little apart, so that symbols with indices do not touch
\AtBeginEnvironment{longtable}{\footnotesize\renewcommand{\arraystretch}{1.45}%
  \let\sptunderscore\_\def\_{\sptunderscore\allowbreak}}
\setlength{\LTcapwidth}{\textwidth}
% A short table is kept on one page.  The build (code/scripts/keep_tables.py) ends the rows of such a table with \\*,
% which forbids a page break after a row, and sets \sptkeptpenalty to 10000 in front of the table.  This penalty is
% put here in the place of the breakpoint (penalty 0) that longtable leaves after every table, so that the paragraph
% that follows a kept table, its caption, begins on the same page; it is set back to 0 after the table.
\newcount\sptkeptpenalty
\sptkeptpenalty=0
\makeatletter
\patchcmd{\endlongtable}{\endgraf\penalty\z@\addvspace\LTpost}%
  {\endgraf\penalty\sptkeptpenalty\addvspace\LTpost\global\sptkeptpenalty=\z@}%
  {}{\PackageError{doc_header}{the end of longtable could not be patched}{The tables can still be typeset, but a short table may be separated from its caption.}}
\makeatother
% allow line breaks after slashes and hyphens inside table cells
\usepackage{xurl}
% code blocks: break over-long lines instead of running into the margin.
% Plain ``` blocks are typeset by pandoc as {verbatim}, highlighted ones
% (```bash) as {Highlighting}; both are re-declared as breakable fancyvrb
% environments.  The continuation line is marked with a small hooked arrow.
\usepackage{fvextra}
\RecustomVerbatimEnvironment{verbatim}{Verbatim}{breaklines,breakanywhere,fontsize=\small}
\makeatletter
\@ifundefined{Highlighting}{}{%
  \RecustomVerbatimEnvironment{Highlighting}{Verbatim}{breaklines,breakanywhere,fontsize=\small,commandchars=\\\{\}}}
\makeatother
% inline code (pandoc writes it as \texttt{...}): never hyphenate -- an inserted
% hyphen could be read as part of a file name -- but allow a line break after
% an underscore, slash, comma, colon or equals sign, so that long paths and
% column lists wrap inside the text width and inside table columns.
% Inside a formula (e.g. the name of a notebook symbol in an equation) the same
% command must first leave math mode, which \text of amsmath does.
\ExplSyntaxOn
\tl_new:N \l__spt_code_tl
\cs_new_protected:Npn \__spt_texttt:n #1
  {
    \tl_set:Nn \l__spt_code_tl {#1}
    \regex_replace_all:nnN { ([\/\,\:\=]) } { \1 \c{allowbreak} } \l__spt_code_tl
    \regex_replace_all:nnN { \c{_} } { \c{_} \c{allowbreak} } \l__spt_code_tl
    \group_begin: \ttfamily \hyphenchar \font = -1 \scan_stop: \l__spt_code_tl \group_end:
  }
\RenewDocumentCommand \texttt { m }
  {
    \mode_if_math:TF { \text { \__spt_texttt:n {#1} } } { \__spt_texttt:n {#1} }
  }
\ExplSyntaxOff
% A figure is written in the Markdown source as an image paragraph followed by
% a caption paragraph ("**Figure S1.** ..."), so that the caption is also
% visible on GitHub.  An image that starts a paragraph is centered.
\usepackage{graphicx}% pandoc loads it only for documents that contain an image
\NewCommandCopy\sptincludegraphics\includegraphics
\RenewDocumentCommand\includegraphics{O{} m}{%
  \ifvmode\begin{center}\sptincludegraphics[#1]{#2}\end{center}%
  \else\sptincludegraphics[#1]{#2}\fi}
% keep figures where they are referenced
\usepackage{float}
\floatplacement{figure}{H}
% running header with the document title
\usepackage{fancyhdr}
\pagestyle{fancy}
\fancyhf{}
\fancyhead[L]{\small PROJECT spacetimes, Revision\_SM\_Higgs}
\fancyhead[R]{\small\thepage}
\renewcommand{\headrulewidth}{0.4pt}
% 16x16 matrices: tighter columns so that a full matrix fits the text width
\setlength{\arraycolsep}{3.2pt}
% The documents of this project contain long tagged symbols, file paths and
% identifiers that cannot be hyphenated.  A paragraph is therefore allowed to
% have looser lines before TeX gives up and lets a line stick out of the margin.
\AtBeginDocument{\tolerance=2000 \emergencystretch=5em \hbadness=2000}
\ifLuaTeX
  \usepackage{selnolig}  % disable illegal ligatures
\fi
\IfFileExists{bookmark.sty}{\usepackage{bookmark}}{\usepackage{hyperref}}
\IfFileExists{xurl.sty}{\usepackage{xurl}}{} % add URL line breaks if available
\urlstyle{same}
\hypersetup{
  pdftitle={The revised hypotheses: a superfluid with the quantum numbers of the Higgs doublet},
  pdfauthor={Prepared for Patrick Lee Nash, Ph.D.},
  colorlinks=true,
  linkcolor={blue},
  filecolor={Maroon},
  citecolor={blue},
  urlcolor={blue},
  pdfcreator={LaTeX via pandoc}}

\title{The revised hypotheses: a superfluid with the quantum numbers of
the Higgs doublet}
\usepackage{etoolbox}
\makeatletter
\providecommand{\subtitle}[1]{% add subtitle to \maketitle
  \apptocmd{\@title}{\par {\large #1 \par}}{}{}
}
\makeatother
\subtitle{Provenance document \texttt{H0} of PROJECT spacetimes, folder
Revision\_SM\_Higgs}
\author{Prepared for Patrick Lee Nash, Ph.D.}
\date{November 21, 2026}

\begin{document}
\maketitle

\hypertarget{what-this-document-is}{%
\section{1. What this document is}\label{what-this-document-is}}

This document states the hypotheses of the folder
\texttt{Revision\_SM\_Higgs} of PROJECT spacetimes and why they are what
they are. The folder was made on 21 November 2026 as a copy of the
folder \texttt{Revision\_HS\_CJT\_and\_Kohn-Sham} at the commit
\texttt{19db425}, which stays unchanged as the record of the earlier
hypotheses. The author asked, at 17:09 PST, that the hypotheses be
modified so that the superfluid can be the Higgs field of the Standard
Model, and he stated the requirement: ``the superfluid must have the
same quantum numbers as the Higgs''. At about 17:15 he chose, from the
constructions that were put to him, the one that this document states.
Everything here is a hypothesis or an exact consequence of the
hypotheses (group theory). The dynamics, whether the superfluid forms
and with which masses, is computed by the programs of this folder from
these hypotheses and is not stated here.

``I'' in this document is the assistant that does the work.

\hypertarget{the-quantum-numbers-of-the-higgs-doublet}{%
\section{2. The quantum numbers of the Higgs
doublet}\label{the-quantum-numbers-of-the-higgs-doublet}}

In the Standard Model the Higgs field \(\Phi^{\rm SM}=(\phi^+,\phi^0)\)
is a Lorentz scalar (spin 0, positive parity), a singlet of colour, a
doublet of \(SU(2)_L\) and has the hypercharge \(Y^{\rm SM}=+\tfrac12\).
With the electric charge \(Q^{\rm SM}=T^{3,{\rm SM}}+Y^{\rm SM}\) its
components have the charges \(+1\) (\(\phi^+\)) and \(0\) (\(\phi^0\)).
It carries neither baryon nor lepton number. With the covariant
derivative
\[D^{\rm SM}_\mu\Phi^{\rm SM}=\big(\partial_\mu+i\,g\,W^a_\mu T^{a,{\rm SM}}+i\,g'\,Y^{\rm SM}B_\mu\big)\Phi^{\rm SM}\tag{2.1}\]
(\(T^{a,{\rm SM}}=\tau^a/2\), \(a=1,2,3\)) and the vacuum
\(\langle\Phi^{\rm SM}\rangle=(0,v/\sqrt2)\), the kinetic term
\((D^{\rm SM}_\mu\Phi^{\rm SM})^\dagger D^{{\rm SM},\mu}\Phi^{\rm SM}\)
gives the masses
\[M_W=\tfrac12g\,v,\qquad M_Z=\tfrac12\sqrt{g^2+g'^2}\,v,\qquad\tan\theta_W=g'/g,\tag{2.2}\]
to the fields \(W^\pm_\mu=(W^1_\mu\mp iW^2_\mu)/\sqrt2\) and
\(Z_\mu=\cos\theta_W\,W^3_\mu-\sin\theta_W\,B_\mu\), and leaves the
photon \(A_\mu=\sin\theta_W\,W^3_\mu+\cos\theta_W\,B_\mu\) massless; the
elementary charge is \(e=g\sin\theta_W\). These are the conventions of
the Review of Particle Physics (S. Navas et al., Particle Data Group,
Physical Review D 110, 030001 (2024), Section 10.1).

\textbf{The measured inputs.} The Review of Particle Physics (2024
edition and its 2025 update) gives the Fermi constant
\(G_F=1.1663788(6)\times10^{-5}\) GeV\(^{-2}\), hence
\(v=(\sqrt2G_F)^{-1/2}=246.22\) GeV; the Higgs mass
\(m_H=125.20\pm0.11\) GeV; \(M_W=80.3692\pm0.0133\) GeV and
\(M_Z=91.1880\pm0.0020\) GeV; and, in the \(\overline{\rm MS}\) scheme
at the scale \(M_Z\), the weak mixing angle
\(\sin^2\hat\theta_W=0.23129(4)\) and the fine-structure constant
\(\hat\alpha(M_Z)^{-1}=127.930(8)\). The couplings \(g=e/\sin\theta_W\)
and \(g'=e/\cos\theta_W\) with \(e^2=4\pi\hat\alpha(M_Z)\) follow from
the last two; the programs of this folder compute them from these
inputs, so that no derived number is typed by hand.

\hypertarget{why-the-superfluid-of-the-earlier-hypotheses-cannot-be-the-higgs}{%
\section{3. Why the superfluid of the earlier hypotheses cannot be the
Higgs}\label{why-the-superfluid-of-the-earlier-hypotheses-cannot-be-the-higgs}}

Under the earlier hypotheses the only spinor field was the isodoublet
\(\Psi_{\rm LL}\), and the superfluid was a condensate of pairs of two
of its particles. The spin-zero, Lorentz-invariant pairs of one Dirac
isodoublet are antisymmetric in the exchange of the two particles (Fermi
statistics). The two Lorentz-invariant matrices of a pair of Dirac
spinors, \(C_{\rm LL}\) and \(C_{\rm LL}\gamma^{\rm ch}_{\rm LL}\), are
antisymmetric, so the isospin factor must be symmetric, and the
symmetric part of the product of two doublets is the triplet:
\[\mathbf 2\otimes\mathbf 2=\mathbf 1_{\rm antisymmetric}\oplus\mathbf 3_{\rm symmetric}.\tag{3.1}\]
The pair channels were the triplets \(\mathsf K^\iota_{\rm LL}\), and
the condensate was an isovector, never a doublet. Moreover the pairs
carried the fermion number 2 and no hypercharge. No choice of couplings
changes this: the quantum numbers of a condensate are fixed by the
fields that pair, not by the strength of the attraction.

\hypertarget{the-revised-hypotheses}{%
\section{4. The revised hypotheses}\label{the-revised-hypotheses}}

\textbf{H1. The fields.} There are two spinor fields, both
four-component Dirac fields of \(SL(2,\mathbb C)\) and both singlets of
colour:

\sptkeptpenalty=10000
\begin{longtable}[]{@{}
  >{\raggedright\arraybackslash}p{(\columnwidth - 10\tabcolsep) * \real{0.1667}}
  >{\raggedright\arraybackslash}p{(\columnwidth - 10\tabcolsep) * \real{0.1667}}
  >{\raggedright\arraybackslash}p{(\columnwidth - 10\tabcolsep) * \real{0.1667}}
  >{\raggedright\arraybackslash}p{(\columnwidth - 10\tabcolsep) * \real{0.1667}}
  >{\raggedright\arraybackslash}p{(\columnwidth - 10\tabcolsep) * \real{0.1667}}
  >{\raggedright\arraybackslash}p{(\columnwidth - 10\tabcolsep) * \real{0.1667}}@{}}
\toprule\noalign{}
\begin{minipage}[b]{\linewidth}\raggedright
field
\end{minipage} & \begin{minipage}[b]{\linewidth}\raggedright
\(SU(2)_L\)
\end{minipage} & \begin{minipage}[b]{\linewidth}\raggedright
hypercharge \(Y^{\rm SM}\)
\end{minipage} & \begin{minipage}[b]{\linewidth}\raggedright
fermion number
\end{minipage} & \begin{minipage}[b]{\linewidth}\raggedright
electric charge \(Q^{\rm SM}\)
\end{minipage} & \begin{minipage}[b]{\linewidth}\raggedright
mass
\end{minipage} \\
\midrule\noalign{}
\endhead
\bottomrule\noalign{}
\endlastfoot
\(\Psi_{\rm LL}\), isospin up & doublet & \(1/4\) & 1 & \(3/4\) &
\(m_{\rm LL}\) \\*
\(\Psi_{\rm LL}\), isospin down & doublet & \(1/4\) & 1 & \(-1/4\) &
\(m_{\rm LL}\) \\*
\(\chi_{\rm SH}\) & singlet & \(1/4\) & 1 & \(1/4\) &
\(m^\chi_{\rm SH}\) \\
\end{longtable}

\(\Psi_{\rm LL}\) is the isodoublet of the earlier hypotheses, with its
mass \(m_{\rm LL}\); \(\chi_{\rm SH}\) is new. Both fields are
vector-like (the two chiralities of each field have the same quantum
numbers), so the gauge group has no anomaly: every anomaly of a
left-handed component is cancelled by the opposite one of its
right-handed partner.

\textbf{H2. The gauge group.} The group of the problem is
\(SL(2,\mathbb C)\times SU(2)_L\times U(1)_Y\). The isospin group
\(SU(2)_{\rm v}\) of the earlier hypotheses is \(SU(2)_L\): its
generators \(T^\iota_{\rm LL}\) on \(\Psi_{\rm LL}\) are the generators
\(T^{a,{\rm SM}}\), its coupling \(g^{\rm v}_{\rm LL}\) is \(g\), and
its gauge field is \(W^a_\mu\). The group \(U(1)^{\rm F}_{\rm LL}\) of
the fermion number \(F_{\rm LL}\) of the earlier hypotheses is the
hypercharge group \(U(1)_Y\), with
\[Y^{\rm SM}=\tfrac14F_{\rm LL}\tag{4.1}\] on both fields; its gauge
field is \(B_\mu\), and the coupling of the earlier hypotheses is
\(g^{\rm F}_{\rm LL}=g'/4\), because
\(g'Y^{\rm SM}=g^{\rm F}_{\rm LL}F_{\rm LL}\). The couplings are the
measured ones of Section 2, no longer free numbers.

\textbf{H3. The interaction.} The pairing term of the earlier
hypotheses, which attracted two particles of \(\Psi_{\rm LL}\) in the
triplet channels \(\mathsf K^\iota_{\rm LL}\), is removed: it would
compete with the doublet and could make a triplet condensate, which does
not have the quantum numbers of the Higgs. In its place the attraction
acts between a particle of \(\Psi_{\rm LL}\) and a particle of
\(\chi_{\rm SH}\) in the spin-zero state built with
\(C_{\rm LL}\gamma^{\rm ch}_{\rm LL}\) (a Lorentz scalar of positive
parity, the channel of the s-wave Cooper pair):
\[\mathcal V^{\rm pair}=-G_{\rm SH}\big[(\Pi^{\rm u}_{\rm SH})^\dagger\Pi^{\rm u}_{\rm SH}+(\Pi^{\rm d}_{\rm SH})^\dagger\Pi^{\rm d}_{\rm SH}\big],\tag{4.2}\]
where \(\Pi^{\rm u}_{\rm SH}\) annihilates an isospin-up particle of
\(\Psi_{\rm LL}\) and a particle of \(\chi_{\rm SH}\) in that state and
\(\Pi^{\rm d}_{\rm SH}\) the same with an isospin-down particle. The sum
is invariant under \(SU(2)_L\times U(1)_Y\) (it is the squared length of
a doublet) and under \(SL(2,\mathbb C)\). The scalar and vector contact
terms of the author's ansatz, with the couplings \(G^{\rm S}_{\rm SF}\)
and \(G^{\rm V}_{\rm SF}\), are kept for \(\Psi_{\rm LL}\). The singlet
field interacts only through (4.2) and the gauge field \(B_\mu\). These
are my decisions of interpretation; the author can change them.

\textbf{H4. The superfluid.} The superfluid is the condensate of the
doublet-singlet pairs. Its pair field is the doublet
\[\Phi_{\rm SH}=\big(\Phi^{\rm u}_{\rm SH},\Phi^{\rm d}_{\rm SH}\big),\qquad\Phi^{\rm u}_{\rm SH}\propto\langle\Pi^{\rm u}_{\rm SH}\rangle,\quad\Phi^{\rm d}_{\rm SH}\propto\langle\Pi^{\rm d}_{\rm SH}\rangle.\tag{4.3}\]
Its quantum numbers follow from H1 and H2 alone:

\begin{itemize}
\tightlist
\item
  spin 0 and positive parity, from the matrix
  \(C_{\rm LL}\gamma^{\rm ch}_{\rm LL}\);
\item
  a colour singlet, because both fields are;
\item
  an \(SU(2)_L\) doublet, because \(\Psi_{\rm LL}\) is a doublet and
  \(\chi_{\rm SH}\) a singlet (\(\mathbf 2\otimes\mathbf 1=\mathbf 2\));
\item
  the hypercharge \(Y^{\rm SM}=\tfrac14+\tfrac14=\tfrac12\): under a
  hypercharge rotation by the angle \(\alpha\) the fields gain the
  phases \(e^{i\alpha/4}\) each, and the pair the phase
  \(e^{i\alpha/2}\), the phase of a field of hypercharge \(\tfrac12\);
\item
  the electric charges \(\tfrac34+\tfrac14=+1\) of
  \(\Phi^{\rm u}_{\rm SH}\) and \(-\tfrac14+\tfrac14=0\) of
  \(\Phi^{\rm d}_{\rm SH}\);
\item
  no baryon and no lepton number, because the fields carry neither; its
  fermion number 2 is not a further quantum number, because by (4.1) it
  is the hypercharge.
\end{itemize}

These are the quantum numbers of \(\Phi^{\rm SM}\), component by
component. The relative number of the two fields (the number of
\(\Psi_{\rm LL}\) particles minus that of \(\chi_{\rm SH}\) particles)
is conserved by the whole Lagrangian and is zero on every pair, so the
condensate leaves it unbroken. When the pair field is not zero it breaks
\(SU(2)_L\times U(1)_Y\) to the \(U(1)\) of \(Q^{\rm SM}\): a rotation
of \(SU(2)_L\times U(1)_Y\) brings it to the form \((0,\varphi)\) with a
real \(\varphi>0\), the form of \(\langle\Phi^{\rm SM}\rangle\), and the
generator that leaves this form unchanged is
\(T^{3,{\rm SM}}+Y^{\rm SM}=Q^{\rm SM}\). The scale of the breaking,
\(v_{\rm SH}\), is defined by the mass of the \(W\) boson,
\(M_W=\tfrac12g\,v_{\rm SH}\), as \(v\) is in the Standard Model.

\textbf{H5. The scale.} The model has one scale, its momentum cutoff
\(\Lambda_{\rm LL}\); the programs compute \(v_{\rm SH}\) in units of
\(\Lambda_{\rm LL}\). The hypothesis that the superfluid is the Higgs
field of the Standard Model fixes the cutoff in GeV by
\(v_{\rm SH}=v=246.22\) GeV. All other masses in GeV, among them the
mass of the amplitude mode of the superfluid, are then predictions, and
the amplitude mode must be compared with \(m_H\).

\hypertarget{consequences-that-follow-at-once}{%
\section{5. Consequences that follow at
once}\label{consequences-that-follow-at-once}}

\begin{enumerate}
\def\labelenumi{\arabic{enumi}.}
\tightlist
\item
  \textbf{The photon and the weak bosons.} If the effective Lagrangian
  of the pair field at second order in the derivatives is the kinetic
  term
  \((D^{\rm SM}_\mu\Phi_{\rm SH})^\dagger D^{{\rm SM},\mu}\Phi_{\rm SH}\)
  with one coefficient, the superfluid gives exactly the masses (2.2),
  with \(v_{\rm SH}\) in place of \(v\), and the massless photon.
  Whether it is of this form is decided by the dynamics: a single
  singlet field, unlike a pair of singlets with equal couplings, does
  not respect the custodial symmetry that makes
  \(\rho=M_W^2/(M_Z^2\cos^2\theta_W)\) equal to 1. In the condensate
  along \(\Phi^{\rm d}_{\rm SH}\) the isospin-down particles of
  \(\Psi_{\rm LL}\) and the particles of \(\chi_{\rm SH}\) are paired,
  the isospin-up particles are not, and this asymmetry contributes to
  \(\rho-1\), as the mass difference of the top and the bottom quark
  does in the Standard Model. The programs compute \(\rho\) from the
  susceptibilities of the superfluid.
\item
  \textbf{Fractional charges.} The constituents carry the electric
  charges \(3/4\), \(-1/4\) and \(1/4\). No particle with such a charge
  has been observed. In this model they are bound in the condensate, but
  its quasi-particles are fermions with these charges and with masses of
  the order of the gap; whether the model is compatible with the
  searches for fractionally charged particles depends on these masses,
  which the dynamics fixes.
\item
  \textbf{No mixing beyond the Standard Model.} Because the fermion
  number is the hypercharge (4.1), the condensate breaks no fourth gauge
  symmetry: the only massive vector fields are \(W^\pm_\mu\) and
  \(Z_\mu\), and there is no extra Goldstone field. The relative number
  of Section 4 (H4) stays a global symmetry.
\item
  \textbf{What does not change.} The bi-octonions, the super-loop, the
  manifold M31 and the Lorentz structure of the spinor fields are those
  of the earlier hypotheses.
\end{enumerate}

\hypertarget{the-symbols-of-this-document}{%
\section{6. The symbols of this
document}\label{the-symbols-of-this-document}}

\sptkeptpenalty=10000
\begin{longtable}[]{@{}
  >{\raggedright\arraybackslash}p{(\columnwidth - 2\tabcolsep) * \real{0.5000}}
  >{\raggedright\arraybackslash}p{(\columnwidth - 2\tabcolsep) * \real{0.5000}}@{}}
\toprule\noalign{}
\begin{minipage}[b]{\linewidth}\raggedright
symbol
\end{minipage} & \begin{minipage}[b]{\linewidth}\raggedright
meaning
\end{minipage} \\
\midrule\noalign{}
\endhead
\bottomrule\noalign{}
\endlastfoot
\(\Psi_{\rm LL}\) & the doublet field: an isodoublet of Dirac fields,
with \(Y^{\rm SM}=1/4\) under these hypotheses \\*
\(\chi_{\rm SH}\) & the singlet field: a Dirac field, a singlet of
\(SU(2)_L\), with \(Y^{\rm SM}=1/4\) \\*
\(m_{\rm LL}\), \(m^\chi_{\rm SH}\) & the masses of \(\Psi_{\rm LL}\)
and of \(\chi_{\rm SH}\) \\*
\(\Lambda_{\rm LL}\) & the momentum cutoff, the unit of the programs \\*
\(C_{\rm LL}\), \(\gamma^{\rm ch}_{\rm LL}\) & the charge-conjugation
matrix and the chirality matrix of the Dirac spinors \\*
\(SL(2,\mathbb C)\) & the double cover of the Lorentz group \\*
\(SU(2)_{\rm v}\), \(U(1)^{\rm F}_{\rm LL}\), \(F_{\rm LL}\) & the
isospin group and the fermion-number group of the earlier hypotheses,
and the fermion number \\*
\(T^\iota_{\rm LL}\), \(\mathsf K^\iota_{\rm LL}\) & the generators of
\(SU(2)_{\rm v}\) on \(\Psi_{\rm LL}\), and the triplet pair matrices of
the earlier hypotheses \\*
\(g^{\rm v}_{\rm LL}\), \(g^{\rm F}_{\rm LL}\) & the gauge couplings of
\(SU(2)_{\rm v}\) and \(U(1)^{\rm F}_{\rm LL}\) \\*
\(G^{\rm S}_{\rm SF}\), \(G^{\rm V}_{\rm SF}\) & the couplings of the
scalar and the vector contact terms of the author's ansatz \\*
\(\Pi^{\rm u}_{\rm SH}\), \(\Pi^{\rm d}_{\rm SH}\) & the pair operators
of the doublet-singlet channel \\*
\(G_{\rm SH}\) & the coupling of the pairing term (4.2) \\*
\(\Phi_{\rm SH}\); \(\Phi^{\rm u}_{\rm SH}\), \(\Phi^{\rm d}_{\rm SH}\)
& the pair field of the superfluid and its components \\*
\(v_{\rm SH}\) & the scale of the breaking by the superfluid,
\(M_W=g\,v_{\rm SH}/2\) \\*
\(SU(2)_L\), \(U(1)_Y\) & the electroweak gauge group of the Standard
Model \\*
\(W^a_\mu\), \(B_\mu\); \(g\), \(g'\); \(\tau^a\), \(T^{a,{\rm SM}}\);
\(Y^{\rm SM}\), \(Q^{\rm SM}\) & the gauge fields, couplings, Pauli
matrices, generators, hypercharge and electric charge of the Standard
Model \\*
\(W^\pm_\mu\), \(Z_\mu\), \(A_\mu\), \(\theta_W\), \(e\), \(v\) & the
mass eigenstates, the weak mixing angle, the elementary charge and the
vacuum value of the Standard Model \\*
\(D^{\rm SM}_\mu\); \(\Phi^{\rm SM}\), \(\phi^+\), \(\phi^0\) & the
covariant derivative of the Higgs doublet, the doublet and its
components \\*
\(M_W\), \(M_Z\), \(m_H\), \(G_F\) & the masses of the \(W\), the \(Z\)
and the Higgs boson, and the Fermi constant \\
\end{longtable}

\end{document}
